% $RUN6 = paper/data/run8-randomization-window
\section{Reproducibility}
\label{sec:repro}
% src: data/CHECKSUMS.json; $RUN6/backtest-upworthy-20261106-2dc461f52ea4.json header.seed 0, header.bootstrap_resamples 500; $RUN6/README.md (backtest seed 0, power seed 20261102); paper/figures/figures.json (every plotted value with its table path, row and field); $RUN6/leakage-audit-upworthy.json, $RUN6/leakage-audit-asos.json, $RUN6/leakage-audit-synthetic.json rungs[].fit_and_scored 0, fit_after_origin 0, outside_window 0, usable_after_mismatch 0, s_star_mismatch 0 for every rung
\iftmlr The code in the anonymized supplementary material\else The released code\fi{} (Python package \texttt{mootcourt}, MIT license) regenerates every results table and every figure from the report files with one command each (\texttt{make tables}, \texttt{make paper-figures}), and the reports themselves with one command per report (Table~\ref{tab:commands}, Appendix~\ref{app:tables}). It lists every plotted value in \texttt{paper/figures/figures.json}. Neither archive is bundled: \texttt{make data} fetches the archive files and refuses any Upworthy file named holdout, confirmatory or undeployed. In a default clone it verifies every file against the committed checksums (Table~\ref{tab:checksums}). With a separate storage root it records the digests on the first download, and the reader compares them with that table. The repository does carry a small sample ledger for tests. Every main backtest and ablation is at seed 0, and Appendix~\ref{app:walltime} lists the identifiers and the other runs' seeds. The leakage audit recomputes that no scored row was visible to the model that scored it and no fit row was usable after its origin: zero violations on all three archives. The audit checks visibility against each test's assigned end, which on the headline archive is the assumed seven days of Section~\ref{sec:setting}, so its zeros hold under that assumption.
% src: .github/workflows/ci.yml (test job on ubuntu-latest, Python 3.11, 3.12 and 3.13; steps after the suite on each version) and scripts/ci_score_report.py: artifact score-report-<python> with the P0, P1 and P2 log score and Brier skill of a fixed-seed backtest of the bundled sample ledger (docs/method.md call 239); tests/test_quick_check.py against tests/fixtures/quick-check/backtest-sample.json (built on macOS; relative 1e-9, 2.0 points on the percentile fields; docs/method.md call 255); docs/troubleshooting.md, Ids differ between machines
On one runner image (ubuntu), continuous integration runs the backtest on the bundled sample ledger under three Python versions and uploads an artifact holding the pooled scores of P0, P1 and P2, so the job itself shows whether the Python versions agree. Within the same job, the quick check gives the cross-platform reading: it rebuilds the sample backtest on the runner and sets it against a fixture made on macOS, at a relative $10^{-9}$. It tolerates 2 points on the percentile fields, so any move smaller than that there passes unseen. Run on a different operating system, the same fit may come out with different last bits in its parameters, and that shifts the model id and the feature hash even though the scores match. So the comparison of cards and models uses the identifiers stored for them, instead of recomputing them (Appendix~\ref{app:walltime}).
% src: docs/method.md call 255 (under the strict comparison, ece_percentile and ece_percentile_independent moved by up to 1.1 points between macOS and ubuntu on the same seed and by 5.2 points in the worst case on the sample report; 26 tied draws of 0.2 points each, a bound and not an observed move; the mid-rank rule counts a draw within 1e-9 * max(1, |observed|) of the observed value as half, src/mootcourt/harness_stats.py null_percentile); scripts/quick_check.py PERCENTILE_TOLERANCE 2.0
We define the calibration percentile as the rank the observed calibration error takes among the draws simulated under its null. Any draw landing within $10^{-9}$ of the observed value is counted as one half. Under the strict comparison this rule replaced, draws like that could land on either side, which shifted the percentile by as much as 1.1 points from macOS to ubuntu with the seed unchanged. For the sample ledger, 5.2 points is the largest shift it could have given a single percentile.

%% tmlr:move-begin artifact
%% tmlr:note artifact How the protocol artifact is kept and how a result is added to it are in Appendix~\ref{sec:artifact}.
%% tmlr:title artifact The protocol artifact and its upkeep
\label{sec:artifact}
\paragraph{Protocol artifact.} Adding a result means a contributor reruns that same backtest against the same snapshot and puts out its report card. Two things go into the artifact, and one stays out.
\begin{itemize}
\item Inside it sit the ledger schema (Table~\ref{tab:schema}), the snapshot digests and admission rules of Section~\ref{sec:setting}, the quarterly origins and the names of the refused files.
% src: README.md, License and citations (MIT)
\item It holds the harness and the report card under the MIT license.
% src: README.md, License and citations (Upworthy Research Archive CC BY 4.0; ASOS Digital Experiments Dataset CC BY 4.0)
\item The raw archive files are absent from it; they are fetched under CC BY 4.0 and never redistributed.
\end{itemize}
The protocol is kept as follows.
\begin{itemize}
\item \iftmlr The authors will maintain it in a public repository after review.\else The author maintains it through the repository, by issues and pull requests.\fi
\item To submit a report card from a new archive, a pull request adds a run folder under \texttt{paper/data} together with its README, in the way the contributing page sets out.
% src: CHANGELOG.md (semantic versioning; before 1.0.0 a schema change is a minor version)
\item \iftmlr The code is under semantic versioning.\else The repository is under semantic versioning. Version 1.0.0 names the released snapshot digests, and before it a schema change counted as a minor version.\fi
\item The snapshot digest names the data a result was read on.
% src: Zenodo record https://zenodo.org/api/records/23988587 (doi 10.5281/zenodo.23988587, conceptdoi 10.5281/zenodo.23988586, read 5 November 2026); $RUN6/README.md (commit 5e4fedf, the code state of this paper's run; commit a7e9f44 for the two learning-curve reports; not repeated: scale/, run-log.json and run-log.md, the timed run, which stay run 7's); $RUN6/README.md (the fourth exception paragraph: five reports at commit 3971829; the fifth: eight files at the analysis code of commit 73d0d70)
\item \iftmlr An earlier version of the code predates\else Release 1.0.0 is archived at Zenodo under DOI\newline \mbox{10.5281/zenodo.23988587}. It predates\fi{} the date rule set out in Section~\ref{sec:setting}, under which tests that began inside the archive's flagged randomization window are excluded. This paper's results rest on \iftmlr one code state\else commit 5e4fedf\fi, apart from six exceptions: one comparison added at \iftmlr a later code state\else commit a7e9f44\fi, the ablation with a Student-t base rate, run at \iftmlr a later code state still\else commit f4566ca\fi, its repeat with exposure fixed, run at \iftmlr a code state that changes no file of the package\else commit e05247e, which changes no file under \texttt{src}\fi, the base rate with a quarter term and the reading by origin, run at \iftmlr a code state after that\else commit 3971829\fi, the runs with exposure fixed on visible rows, the interval with quarters as units and the no-effect reference, run at \iftmlr the latest code state\else commit 73d0d70\fi, and the wall times and scale measurements of Appendix~\ref{app:walltime}, taken from earlier code states that the appendix describes. The run folder \texttt{paper/data/run8-randomization-window} holds the reports.
% src: Zenodo record https://zenodo.org/api/records/23988587 (conceptdoi 10.5281/zenodo.23988586, read 5 November 2026)
\item \iftmlr The code, the reports and the plan are in the anonymized supplementary material.\else The concept DOI,\newline \mbox{10.5281/zenodo.23988586}, resolves to the latest release.\fi
% src: references.bib mootcourt2026plan (posted 2026-11-04)
\item \iftmlr The plan is the copy in that material.\else The plan is the copy posted publicly on 4 November 2026 \cite{mootcourt2026plan}.\fi
\end{itemize}
%% tmlr:move-end artifact
